\section{Discussão}

A síntese dos seis temas permite quatro leituras transversais que atravessam as perguntas de pesquisa e organizam o que o corpus, tomado como um todo, informa sobre o estado do campo.

\noindent\textbf{A unidade de trabalho deslocou-se do prompt para o processo.} A tese central deste estudo sustenta-se empiricamente. Em 34 dos 37 estudos a contribuição é descrita como framework, e os campos de artefatos persistentes, papéis e validação estão presentes na esmagadora maioria das fichas. O modelo de linguagem raramente é o diferenciador declarado; o que distingue as propostas é a estrutura de processo que organiza specs, papéis, execução e validação em torno do trabalho do agente. Os três estudos que não são frameworks reforçam o contraste: uma arquitetura de padrões de prompt, uma plataforma de cadeias de IA e um método de cenários futuros. Esses são exatamente os casos em que a unidade ainda é a técnica isolada e não o processo \cite{JulesWhite2023,YuCheng2023,JSauvola2024}.

\noindent\textbf{A taxonomia de seis dimensões é parcialmente validada e parcialmente desafiada.} Cinco dimensões (especificação, contexto, papéis, execução e validação) têm ancoragem empírica forte em pelo menos um tema. A sexta, portabilidade, é fraca e quase sempre diagnóstica: aparece como risco de aprisionamento e de dependência de plataforma, não como propriedade de projeto deliberada. Além disso, dois eixos relevantes emergem fora das seis dimensões, o humano-organizacional, em T5, e o de segurança, em T6. A recomendação analítica que decorre do modo híbrido é concreta: considerar estender a taxonomia com uma dimensão de colaboração e governança humano-IA e tratar segurança como critério de avaliação transversal ou subdimensão de validação. Essa é a contribuição que a confrontação dos temas com a taxonomia agrega em relação a assumir as dimensões a priori.

\noindent\textbf{Três eixos de tensão recorrentes.} O corpus exibe consenso parcial, não convergência, ao longo de três eixos. O primeiro é autonomia versus controle humano: CodePori, MetaGPT e ChatDev inclinam-se ao fluxo autônomo de ponta a ponta, enquanto frameworks centrados em adoção e confiança inserem a decisão humana como elemento de projeto, e a distinção entre codificar por intuição e codificar de forma agêntica teoriza o contraste \cite{ZeeshanRasheed2024,SiruiHong2023,ChenQian2023,SaiZhang2025,DanielRusso2024,Baron2025,RanjanSapkota2025}. O segundo é governança pesada versus agilidade leve: cerimônias de verificação e contratos formais contrapõem-se a fluxos enxutos para startups e a padrões de prompt, e o custo de coordenação é o preço recorrente da governança \cite{Merchant2025,Paduraru2026,Khan2026,JulesWhite2023}. O terceiro é a direção do fluxo da especificação: especificar antes, em desenvolvimento orientado por especificação, MetaGPT e AgileGen, contra recuperar a especificação a partir do legado, em Reversa \cite{DeepakBabuPiskala2026,SiruiHong2023,SaiZhang2025,Macedo2026}. Os dois extremos ancoram-se na mesma dimensão de especificação, mas em sentidos opostos, o que é material para qualquer matriz comparativa futura.

\noindent\textbf{Rastreabilidade e contexto são os requisitos transversais mais citados.} O risco associado ao contexto aparece em quase todos os estudos, e a rastreabilidade é demanda recorrente nos temas de especificação e de governança. Juntos, indicam que os gargalos práticos do campo não estão na geração de código em si, mas em manter o agente ancorado no projeto real e auditável ao longo do tempo. Esse achado realinha a leitura usual do problema: o desafio aberto é menos produzir código plausível e mais sustentar fidelidade ao sistema e à intenção registrada.

\noindent\textbf{A evidência é desigual.} Embora 26 dos 37 estudos declarem avaliação empírica e 22 tenham repositório público, dez são preprints arXiv ainda não revisados por pares e a maior parte das avaliações ocorre em benchmarks ou estudos de caso pontuais, não em projetos industriais longitudinais. A força da evidência varia muito por categoria de fonte, exatamente como o protocolo antecipou ao classificar a evidência por categoria em vez de aplicar exclusão por qualidade. Essa heterogeneidade recomenda cautela: os achados deste estudo descrevem com segurança a estrutura e as tensões do campo, mas afirmações sobre desempenho relativo entre frameworks permanecem subdeterminadas pela evidência disponível.
